% ============================================================================
% Fichier  : partie3/P03_C08_exploitation.tex
% Titre    : L'Exploitation : Principes et Cadres
% Auteur   : MINKA MI NGUIDJOÏ Thierry Emmanuel
% Version  : 1.0 -- novembre 2026
% ============================================================================
% Sources : NIST SP 800-115 (attaque) ; PTES (exploitation) ; MITRE ATT&CK.
% Démonstration de l'Opération LOGONE : cas fictif, aucun code d'attaque.
% ============================================================================

\chapter{L'Exploitation : Principes et Cadres}
\label{chap:exploitation}
\chaptermark{L'exploitation}

\epigraph{%
    « Démontrer, pas détruire. »%
}{Principe du cours}

\niveauChapitre{Intermédiaire}{Chapitres~\ref{chap:engagement}
et~\ref{chap:vulnerabilites}.}

\begin{boxobjectifs}
\begin{itemize}
    \item Distinguer vulnérabilité, exploit, charge utile et preuve de
          concept.
    \item Choisir de démontrer ou de ne pas démontrer une faille.
    \item Évaluer le risque d'une exploitation pour la production.
    \item Utiliser un cadre d'exploitation de manière maîtrisée.
    \item Produire une preuve minimale, vérifiable et sans dommage.
\end{itemize}
\end{boxobjectifs}

\bigskip

% ============================================================================
\section*{Question directrice : faut-il exploiter l'ancien portail pour prouver qu'il est vulnérable~?}
% ============================================================================

L'ancien portail de suivi, oublié depuis 2024, utilise une version
d'application connue pour une faille d'exécution de code à distance. Le
scanner le dit. M.~KOUNOU répond que « ce serveur ne sert plus à rien ».
M.~BIKORO décide de démontrer la faille, avec la plus petite preuve
possible. Pourquoi, et comment~?

\medskip
\textbf{Plan du chapitre.}
\begin{enumerate}
    \item Le vocabulaire de l'exploitation (section~\ref{sec:c08-vocabulaire}).
    \item Démontrer ou non (section~\ref{sec:c08-decider}).
    \item Le risque pour la production (section~\ref{sec:c08-risque}).
    \item Les cadres d'exploitation (section~\ref{sec:c08-cadres}).
    \item La preuve minimale (section~\ref{sec:c08-preuve}).
\end{enumerate}

\separateur

% ============================================================================
\section{Le Vocabulaire de l'Exploitation}
\label{sec:c08-vocabulaire}
% ============================================================================

\begin{boxdef}
\termecle{Exploit}~: moyen (programme, requête, séquence d'actions) qui
tire parti d'une vulnérabilité pour obtenir un effet non prévu.
\termecle{Charge utile}~: ce que l'exploit exécute une fois la faille
déclenchée. \termecle{Preuve de concept}~: démonstration minimale qu'une
faille est exploitable, sans en tirer tout le parti possible.
\end{boxdef}

\begin{tableaucours}{L{3.4cm}L{8.8cm}}{%
    \tableauHeader{Effet recherché} &
    \tableauHeader{Exemple de preuve minimale}
}
Exécution de code & Afficher l'identité du compte d'exécution et le nom de
la machine \\
Lecture de fichier & Lire un fichier sans valeur (version, fichier de test
déposé par accord) \\
Contournement d'authentification & Afficher la page d'accueil d'un compte
de test fourni \\
Élévation de privilèges & Afficher l'appartenance au groupe obtenu \\
\end{tableaucours}
\vspace{-0.8em}
\begin{center}{\small\itshape Tableau~8.1. Effets et preuves minimales}\end{center}

% ============================================================================
\section{Démontrer ou Non}
\label{sec:c08-decider}
% ============================================================================

Démontrer une faille la rend indiscutable et en mesure l'impact réel~; ne
pas la démontrer évite tout risque pour la production. La décision dépend
de l'objectif, des règles d'engagement et du risque.

\begin{boxprocedure}
\textbf{Décider d'exploiter}
\begin{enumerate}
    \item L'exploitation est-elle autorisée pour cette cible~?
    \item L'exploit est-il connu, stable, compris~? Sait-on ce qu'il fait~?
    \item Quel est le pire effet possible (arrêt, corruption)~? Est-il
          acceptable~?
    \item Le client est-il prévenu, la fenêtre adaptée, une sauvegarde
          faite~?
    \item La preuve minimale suffit-elle à l'objectif~?
\end{enumerate}
\end{boxprocedure}

\begin{boxalert}
\textbf{Violation typique.} Lancer un exploit téléchargé sans l'avoir lu.

\smallskip
\textbf{Règle~:} un exploit d'origine inconnue peut contenir une porte
dérobée, détruire des données ou faire tomber le service. On n'emploie sur
la production que du code dont on connaît la provenance et l'effet, après
l'avoir essayé dans un laboratoire (Chapitre~\ref{chap:labo}).
\end{boxalert}

% ============================================================================
\section{Le Risque pour la Production}
\label{sec:c08-risque}
% ============================================================================

\begin{tableaucours}{L{3.4cm}L{4.6cm}L{4.2cm}}{%
    \tableauHeader{Type de faille} &
    \tableauHeader{Risque de l'exploitation} &
    \tableauHeader{Précaution}
}
Dépassement de mémoire & Arrêt du service, instabilité & Laboratoire
d'abord~; fenêtre creuse \\
Injection dans une base & Modification de données & Requêtes en lecture
seule \\
Exécution de commande & Commande destructrice par erreur & Commandes
d'identification seulement \\
Déni de service & Interruption certaine & Exclu par défaut \\
\end{tableaucours}
\vspace{-0.8em}
\begin{center}{\small\itshape Tableau~8.2. Risques et précautions}\end{center}

% ============================================================================
\section{Les Cadres d'Exploitation}
\label{sec:c08-cadres}
% ============================================================================

Les \termecle{cadres d'exploitation} réunissent des modules d'exploits, de
charges utiles et d'outils de post-exploitation, documentés et maintenus.
Ils réduisent le risque des exploits isolés, mais ne dispensent ni de
comprendre le module, ni de vérifier qu'il est adapté à la version ciblée.
Le cours ne détaille pas leur emploi~: il en retient l'usage
professionnel~: modules connus, options vérifiées, journal de chaque
lancement.

% ============================================================================
\section{La Preuve Minimale}
\label{sec:c08-preuve}
% ============================================================================

\begin{boxscenario}{Opération LOGONE. L'ancien portail}
\textbf{Constatation.} Le 22~octobre à 14~h~05, après accord de M.~KOUNOU et
sauvegarde du serveur, M.~TAMO déclenche la faille de l'ancien portail avec
un module connu. La preuve~: l'affichage de l'identité du compte de service
web et du nom de la machine. Aucune autre commande n'est exécutée~; la
session est fermée à 14~h~08.

\medskip
\textbf{Analyse} (confiance élevée). « Ce serveur ne sert plus à rien »~:
mais il est en ligne, dans la DMZ, et le compte web peut lire la
configuration de connexion à la base du portail actuel. La démonstration a
transformé une alerte du scanner en constat incontestable.
\end{boxscenario}

\begin{boxobjet}{T-09 · Exécution de code sur l'ancien portail de suivi}
\textbf{Description~:} « ancien-portail.lizbiz.cm » exécute une version
d'application vulnérable à une exécution de code à distance.
\textbf{Preuve~:} capture du 22~octobre, 14~h~06 (identité du compte, nom
de la machine). \textbf{Gravité~:} élevée dans le contexte (serveur exposé).
\textbf{Recommandation~:} retirer le serveur~; à défaut, le mettre à jour et
l'isoler. \textbf{Détection~:} processus inhabituel lancé par le compte web.
\end{boxobjet}

\begin{boxmetier}{Pentesteur}
\textbf{Une commande, une preuve, on ferme.} Le pentesteur prépare à
l'avance la commande de preuve, l'exécute une fois, capture, ferme la
session et note l'heure. Explorer « tant qu'on y est » n'est pas une preuve,
c'est un risque.
\end{boxmetier}

\separateur

\begin{boxresume}
\textbf{Ce qu'il faut retenir de ce chapitre~:}
\begin{itemize}
    \item \textbf{Vulnérabilité, exploit, charge utile, preuve de concept}~:
          quatre notions distinctes.
    \item \textbf{On décide d'exploiter} selon l'autorisation, la stabilité,
          le pire effet et l'objectif.
    \item \textbf{Jamais d'exploit inconnu} sur la production~; laboratoire
          d'abord.
    \item \textbf{Les cadres d'exploitation} réduisent le risque sans
          dispenser de comprendre.
    \item \textbf{Une preuve minimale} suffit~: une commande, une capture,
          on ferme.
\end{itemize}
\end{boxresume}

\begin{boxinfo}
\textbf{Points de contrôle de l'élève.} Je sais~: définir exploit, charge
utile et preuve de concept~; dérouler la décision d'exploiter~; associer
une précaution à un type de faille~; choisir une preuve minimale.
\end{boxinfo}

\bigskip

\begin{boxexo}[title={\ding{45}\quad Exercice~8.1. Vocabulaire \niveaubase}]
Dans l'épisode de l'ancien portail, identifiez la vulnérabilité, l'exploit,
la charge utile et la preuve.
\end{boxexo}

\begin{boxexo}[title={\ding{45}\quad Exercice~8.2. Preuve minimale \niveaubase}]
Proposez une preuve minimale pour une lecture de fichier arbitraire sur un
serveur web.
\end{boxexo}

\begin{boxexo}[title={\ding{45}\quad Exercice~8.3. Déni de service \niveaubase}]
Pourquoi le déni de service est-il exclu par défaut~? Quel article de la loi
n°~2010/012 le concerne~?
\end{boxexo}

\begin{boxexo}[title={\ding{45}\quad Exercice~8.4. Décider \niveaumoyen}]
Appliquez la procédure de décision à une faille de dépassement de mémoire
sur le serveur de base de données de la facturation.
\end{boxexo}

\begin{boxexo}[title={\ding{45}\quad Exercice~8.5. L'exploit inconnu \niveaumoyen}]
Citez trois dangers d'un exploit téléchargé sur un site non vérifié.
\end{boxexo}

\begin{boxexo}[title={\ding{45}\quad Exercice~8.6. « Ne sert plus à rien » \niveaumoyen}]
Comment répondre à M.~KOUNOU, qui juge l'ancien portail sans importance~?
\end{boxexo}

\begin{boxexo}[title={\ding{45}\quad Exercice~8.7. La détection \niveaumoyen}]
Quels journaux auraient permis de détecter l'exploitation de l'ancien
portail~?
\end{boxexo}

\begin{boxexo}[title={\ding{45}\quad Exercice~8.8. Ne pas démontrer \niveaumoyen}]
Donnez deux situations où le testeur doit renoncer à démontrer une faille,
et comment il la rapporte alors.
\end{boxexo}

\begin{boxexo}[title={\ding{45}\quad Exercice~8.9. Une procédure \niveauexpert}]
Rédigez la procédure d'exploitation contrôlée de NYONG SÉCURITÉ~: accord,
préparation, exécution, preuve, fermeture, journal.
\end{boxexo}

\begin{boxexo}[title={\ding{45}\quad Exercice~8.10. Laboratoire \niveauexpert}]
Décrivez le laboratoire qui permet de valider un module d'exploitation avant
son emploi chez un client.
\end{boxexo}

\bigskip

\begin{boxinfo}
\textbf{Pour aller plus loin.}
\begin{itemize}
    \item NIST SP~800-115~\cite{nist800115}~; PTES~\cite{ptes}.
\end{itemize}
\end{boxinfo}

\begin{boxinfo}
\textbf{Transition.} Sur le réseau interne, l'attaquant cherche des secrets
qui circulent~: le Chapitre~\ref{chap:reseau} traite des attaques réseau et
de l'authentification.
\end{boxinfo}
